% GENERATED FILE. Edit canonical lessons, then run npm run generate:books.
% canonical-source-hash: fnv1a64:5ef5fc4e50a7f703
% canonical-lessons: IT-C317-spumante, IT-C317-tappo, IT-C317-stappare, IT-C317-cavatappi, IT-C317-banchetto

\chapter{Cin cin}
\label{ch:it-cin-cin}
\hlchaptermodality{317}

\begin{chapteropening}
\textbf{By the end of this chapter:} \emph{I can describe opening the bubbly: the sparkling wine, the cork, uncorking, the corkscrew and the feast.}

Complete the last lesson of chapter 317: I can describe opening the bubbly: the sparkling wine, the cork, uncorking, the corkscrew and the feast.
\end{chapteropening}

\section[lo spumante]{lo spumante — the sparkling wine}

\label{lesson:IT-C317-spumante}



\begin{quote}
Before the new one: say the Italian for the bow, then the Italian for the balloon.
\end{quote}

\subsection*{You'll want to know: lo spumante}
\textbf{lo spumante} — \textquotedblleft{}the sparkling wine\textquotedblright{}.

\begin{quote}
A mezzanotte apriamo lo spumante.
\end{quote}

\begin{quote}
\emph{At midnight we open the sparkling wine.}
\end{quote}

Built on \emph{spuma}, ‘foam’.

\begin{cousinweb}
\emph{spumante} is the present participle of \emph{spumare} \textquotedblleft{}to foam\textquotedblright{}, from Latin \emph{spūma} \textquotedblleft{}foam\textquotedblright{} (Lewis \& Short).
\end{cousinweb}

\begin{grammarlens}[title={What you've built}]
One more word for telling, explaining and describing at B1.
\end{grammarlens}

\subsection*{Your turn}
\begin{itemize}
\raggedright
  \item \emph{Say it:} \emph{lo spumante}
  \item \emph{Say it:} \emph{lo spumante}, once more
  \item \emph{Say it:} a sentence of your own with \emph{lo spumante}
  \item \emph{Recall:} say the Italian for the bow, then the Italian for the balloon, then say \emph{lo spumante} again
\end{itemize}

\begin{tcolorbox}[breakable,colback=teal!4,colframe=teal!35!black,title={Before you move on}]
What does lo spumante mean? (\textquotedblleft{}The sparkling wine\textquotedblright{}.) Say it once more.
\end{tcolorbox}
\section[il tappo]{il tappo — the cork, the cap}

\label{lesson:IT-C317-tappo}



\begin{quote}
Before the new one: say the Italian for the balloon, then the Italian for the sparkling wine.
\end{quote}

\subsection*{You'll want to know: il tappo}
\textbf{il tappo} — \textquotedblleft{}the cork, the cap\textquotedblright{}.

\begin{quote}
Il tappo è volato fino al soffitto.
\end{quote}

\begin{quote}
\emph{The cork flew up to the ceiling.}
\end{quote}

Informally, \emph{un tappo} is also a short person.

\begin{cousinweb}
\emph{tappo} comes from a Frankish \emph{tappo} \textquotedblleft{}stopper, plug\textquotedblright{} (EtymDB).
\end{cousinweb}

\begin{grammarlens}[title={What you've built}]
One more word for telling, explaining and describing at B1.
\end{grammarlens}

\subsection*{Your turn}
\begin{itemize}
\raggedright
  \item \emph{Say it:} \emph{il tappo}
  \item \emph{Say it:} \emph{il tappo}, once more
  \item \emph{Say it:} a sentence of your own with \emph{il tappo}
  \item \emph{Recall:} say the Italian for the balloon, then the Italian for the sparkling wine, then say \emph{il tappo} again
\end{itemize}

\begin{tcolorbox}[breakable,colback=teal!4,colframe=teal!35!black,title={Before you move on}]
What does il tappo mean? (\textquotedblleft{}The cork, the cap\textquotedblright{}.) Say it once more.
\end{tcolorbox}
\section[stappare]{stappare — to uncork, to open (a bottle)}

\label{lesson:IT-C317-stappare}



\begin{quote}
Before the new one: say the Italian for the sparkling wine, then the Italian for the cork.
\end{quote}

\subsection*{You'll want to know: stappare}
\textbf{stappare} — \textquotedblleft{}to uncork, to open (a bottle)\textquotedblright{}.

\begin{quote}
Chi stappa la bottiglia?
\end{quote}

\begin{quote}
\emph{Who is opening the bottle?}
\end{quote}

The opposite of \emph{tappare}, to stop up.

\begin{cousinweb}
\emph{stappare} is \emph{s-} \textquotedblleft{}un-\textquotedblright{} with \emph{tappo} \textquotedblleft{}cork\textquotedblright{}, from Frankish \emph{tappo} (EtymDB).
\end{cousinweb}

\begin{grammarlens}[title={What you've built}]
One more word for telling, explaining and describing at B1.
\end{grammarlens}

\subsection*{Your turn}
\begin{itemize}
\raggedright
  \item \emph{Say it:} \emph{stappare}
  \item \emph{Say it:} \emph{stappare}, once more
  \item \emph{Say it:} a sentence of your own with \emph{stappare}
  \item \emph{Recall:} say the Italian for the sparkling wine, then the Italian for the cork, then say \emph{stappare} again
\end{itemize}

\begin{tcolorbox}[breakable,colback=teal!4,colframe=teal!35!black,title={Before you move on}]
What does stappare mean? (\textquotedblleft{}To uncork, to open (a bottle)\textquotedblright{}.) Say it once more.
\end{tcolorbox}
\section[il cavatappi]{il cavatappi — the corkscrew}

\label{lesson:IT-C317-cavatappi}



\begin{quote}
Before the new one: say the Italian for the cork, then the Italian for to uncork.
\end{quote}

\subsection*{You'll want to know: il cavatappi}
\textbf{il cavatappi} — \textquotedblleft{}the corkscrew\textquotedblright{}.

\begin{quote}
Non trovo il cavatappi!
\end{quote}

\begin{quote}
\emph{I cannot find the corkscrew!}
\end{quote}

\emph{Cava} (pulls out) + \emph{tappi} (corks). Plural unchanged.

\begin{cousinweb}
\emph{cavatappi} joins \emph{cava}, from \emph{cavare} \textquotedblleft{}to pull out\textquotedblright{}, Latin \emph{cavāre} (Lewis \& Short), and \emph{tappi} \textquotedblleft{}corks\textquotedblright{}.
\end{cousinweb}

\begin{grammarlens}[title={What you've built}]
One more word for telling, explaining and describing at B1.
\end{grammarlens}

\subsection*{Your turn}
\begin{itemize}
\raggedright
  \item \emph{Say it:} \emph{il cavatappi}
  \item \emph{Say it:} \emph{il cavatappi}, once more
  \item \emph{Say it:} a sentence of your own with \emph{il cavatappi}
  \item \emph{Recall:} say the Italian for the cork, then the Italian for to uncork, then say \emph{il cavatappi} again
\end{itemize}

\begin{tcolorbox}[breakable,colback=teal!4,colframe=teal!35!black,title={Before you move on}]
What does il cavatappi mean? (\textquotedblleft{}The corkscrew\textquotedblright{}.) Say it once more.
\end{tcolorbox}
\section[il banchetto]{il banchetto — the banquet, the feast}

\label{lesson:IT-C317-banchetto}



\begin{quote}
Before the new one: say the Italian for to uncork, then the Italian for the corkscrew.
\end{quote}

\subsection*{You'll want to know: il banchetto}
\textbf{il banchetto} — \textquotedblleft{}the banquet, the feast\textquotedblright{}.

\begin{quote}
Il banchetto è durato fino a notte.
\end{quote}

\begin{quote}
\emph{The banquet went on into the night.}
\end{quote}

A little \emph{banco}, ‘bench, table’; also a market stall.

\begin{cousinweb}
\emph{banchetto} is a little \emph{banco} \textquotedblleft{}bench\textquotedblright{} (EtymDB), which comes from a Longobard word for bench.
\end{cousinweb}

\begin{grammarlens}[title={What you've built}]
One more word for telling, explaining and describing at B1.
\end{grammarlens}

\subsection*{Your turn}
\begin{itemize}
\raggedright
  \item \emph{Say it:} \emph{il banchetto}
  \item \emph{Say it:} \emph{il banchetto}, once more
  \item \emph{Say it:} a sentence of your own with \emph{il banchetto}
  \item \emph{Recall:} say the Italian for to uncork, then the Italian for the corkscrew, then say \emph{il banchetto} again
\end{itemize}

\begin{tcolorbox}[breakable,colback=teal!4,colframe=teal!35!black,title={Before you move on}]
What does il banchetto mean? (\textquotedblleft{}The banquet, the feast\textquotedblright{}.) Say it once more.
\end{tcolorbox}
